\begin{recipe}{Simpele kippensoep}
\kicker[Voorgerecht,Soep,Vlees,Nederlands]{Voorgerecht · soep · vlees · Nederlands}
\index[register]{Simpele kippensoep}
\index[register]{Kipdijfilet}
\index[register]{Kippenfond}
\index[register]{Soepgroente}

\meta{20 minuten \dvd Eén pan}

\lettrine{D}{eze} kippensoep zetten we vaak op voordat we pannenkoeken bakken, zodat we
niet alleen maar zoet eten. Snel gemaakt, met weinig ingrediënten.

\blockrule{Ingrediënten}
\begin{ingredients}
  \ing{200 g}{kipdijfilet, in kleine stukjes}\index[register]{Kipdijfilet}
  \ingb{bakboter}
  \ing{250 g}{soepgroente}\index[register]{Soepgroente}
  \ing{400 ml}{kippenfond}\index[register]{Kippenfond}
  \ingb{2 à 3 lege potten water (gebruik de hals van de fondpot als maat)}
  \ingb{vermicelli, optioneel}
\end{ingredients}

\blockrule{Bereiding}
\tip{Koop je een grotere zak soepgroente, dan vries je de rest prima in voor de
volgende keer.}
\begin{steps}
  \item Snijd de kipdijfilet in kleine stukjes en bak ze in bakboter in een soeppan
        bruin aan. Het aanbaksel op de bodem geeft straks smaak aan de soep.
  \item Voeg de soepgroente toe en bak kort mee.
  \item Schenk de kippenfond erbij, en 2 à 3 potten water (gebruik de hals van de
        fondpot als maat).
  \item Doe het deksel erop en laat 10 minuten zachtjes doorpruttelen, niet hard
        koken.
  \item Roer er eventueel wat vermicelli door voor meer body en laat die
        meegaren.
\end{steps}
\end{recipe}
